\section{Plano de diseño de un sistema de Agent: del concepto de investigación al producto ejecutable}

Las secciones anteriores trataron por separado el lenguaje, el entorno de la tarea, el reward y el multi-agent; esta sección los integra en un plano de sistema que puede diseñarse. Un producto de Agent no suele ser un solo prompt, sino un ciclo: el objetivo del usuario entra al sistema, el sistema convierte ese objetivo en un estado de la tarea, el modelo genera un plan, las herramientas ejecutan acciones, el entorno devuelve observaciones, el verificador evalúa el resultado, la capa de memoria guarda la experiencia y, por último, se pasa a la siguiente ronda.

\begin{importantbox}{Ciclo cerrado mínimo de un sistema de Agent}
\[
\text{Goal} \rightarrow \text{State} \rightarrow \text{Plan} \rightarrow \text{Action} \rightarrow \text{Observation} \rightarrow \text{Reward/Verification} \rightarrow \text{Memory}
\]
Aquí, Goal es el objetivo del usuario, State es el estado actual de la tarea, Plan es el plan que genera el modelo, Action es la ejecución de herramientas o de código, Observation es la retroalimentación del entorno, Reward/Verification es la señal de evaluación y Memory es la experiencia a largo plazo junto con el contexto.
\end{importantbox}

\begin{knowledgebox}{Plano de diseño: qué pregunta cada capa}
\begin{tabularx}{\textwidth}{p{0.20\textwidth}p{0.34\textwidth}X}
\toprule
Capa & Pregunta de diseño & Fallo frecuente \\
\midrule
Capa de objetivo & ¿Qué quiere realmente el usuario? ¿Cuál es la condición de éxito? & El objetivo es ambiguo y el Agent hace algo que parece útil pero no tiene relación con lo pedido. \\
Capa de estado & ¿Qué archivos, herramientas, permisos y restricciones hay en este momento? & Falta contexto y las acciones se basan en un estado erróneo. \\
Capa de planificación & ¿Hace falta dividir la tarea, trabajar en paralelo o consultar al usuario? & El plan es demasiado largo y frágil, y no se ajusta ante una excepción. \\
Capa de acción & ¿Qué acciones pueden ejecutarse con seguridad y cuáles requieren confirmación? & Exceso de permisos, borrados por error, llamadas a la herramienta equivocada. \\
Capa de verificación & ¿Cómo se sabe que la tarea realmente se completó? & Solo se revisa la salida superficial, sin verificar el resultado real. \\
Capa de memoria & ¿Qué experiencia debe conservarse para la próxima vez? & Memoria contaminada, riesgos de privacidad, información desactualizada. \\
\bottomrule
\end{tabularx}
\end{knowledgebox}

\begin{warningbox}{La distancia entre una demo y un producto}
Una demo de Agent puede impresionar con una sola trayectoria exitosa; un producto de Agent tiene que manejar la recuperación ante fallos, los permisos, los registros, la verificación, las interrupciones del usuario, la memoria a largo plazo y el costo. Buena parte de las capacidades de sistema de la segunda mitad reside precisamente en estas capas «aburridas pero necesarias».
\end{warningbox}

\subsection{Resumen de la sección}

Lo esencial del diseño de un sistema de Agent no es envolver el modelo con una UI, sino cerrar el ciclo entre objetivo, estado, plan, acción, observación, verificación y memoria. Solo cuando el ciclo es estable puede el Agent pasar de ser una demostración deslumbrante a un producto confiable.
